\chapter{The arena: geometry, draw, and ground loss}
\label{ch:geometry}

\section{Track geometry and the movable rail}

The draw bias is fundamentally \emph{geometric}: on a turn, an inner lane is a shorter path, so a
low-drawn horse that holds the rail runs less ground. The model carries an explicit geometry,
parameterised by \emph{venue $\times$ surface $\times$ course}, because HKJC turf tracks use a
\emph{movable running rail} --- the race's \code{Course} (A, A+2, B, B+2, B+3, C, C+3) shifts the
rail outward to spread wear, changing the usable width and slightly the home-straight length. Three
structural regimes are encoded as \code{TrackGeometry} objects:
\begin{itemize}
  \item \textbf{Happy Valley} (turf, $\SI{1450}{m}$ circumference, right-handed): a tight,
        bend-dominated circuit ($\approx 470 + 530\,\si{m}$ of turn within a $\SI{1450}{m}$ lap)
        --- the strongest draw bias in Hong Kong.
  \item \textbf{Sha Tin turf} (galloping oval, $\SI{1900}{m}$, right-handed, $\SI{430}{m}$ home
        straight): long straights and two $\approx \SI{495}{m}$ bends, hence a \emph{weaker}
        geometric draw bias than HV.
  \item \textbf{Sha Tin straight $\SI{1000}{m}$ chute}: a separate course with \emph{no bends}, so
        no bend-geometry draw bias at all (only going/watering gradients, deliberately unmodelled).
  \item \textbf{Sha Tin all-weather} (dirt, $\SI{1560}{m}$, fixed single rail, uniform
        $\SI{22.8}{m}$ width).
\end{itemize}
The geometry exposes a \emph{bend fraction} --- the fraction of a race of length $L$ run on turns,
computed by sampling positions along the trip and testing each against the bend segments,
\begin{equation}
  \mathrm{bend\_fraction}(L) \;=\; \frac{1}{P}\sum_{p=1}^{P}\mathbb{1}\bigl[\text{position } d_p
   \text{ is on a turn}\bigr],
  \label{eq:bendfrac}
\end{equation}
which is the draw-bias exposure of the race. It is $\approx 0$ on the straight chute (which is why
the draw term is zeroed there) and largest at HV.

\section{Draw as ground loss on the drift}

A horse drawn $\mathrm{Dr}$ positions off the rail runs a longer arc round the turns. The
production model represents this as an \emph{effective lengthening} of the trip the drift must
cover,
\begin{equation}
  \Leff \;=\; L\,\bigl(1 + g_{\text{draw}}\,(\mathrm{Dr}-1)\bigr),
  \label{eq:leff}
\end{equation}
which feeds straight into the Wald mean \eqref{eq:mu-wprime}: a wider draw means a longer effective
trip, a longer first-passage time, and a lower win probability. The gain $g_{\text{draw}}$ is
\emph{per-venue and geometric}, baked as $\texttt{DRAW\_GAIN\_BY\_VENUE}$ (ST $0.003$, HV $0.008$):
HV, bend-dominated, wants roughly $2.7\times$ the ST penalty. On the all-weather surface the rail
bias is stronger and is overridden by $\texttt{DRAW\_GAIN\_AWT}=0.008$. On any course with
$\mathrm{bend\_fraction}<0.05$ (the straight chute) the gain is set to zero --- no arc, no penalty.

\begin{remark}[A lesson in always sweeping]
An early diagnosis read the HV draw residual as ``over-weighted'' and proposed setting the gain to
zero; the sweep proved the \emph{opposite} --- zeroing it craters HV to $41\%$, because HV was in
fact massively \emph{under}-penalised by the old global $0.003$. The episode is the canonical
reminder that the sign of a residual coupling must be established by a sweep against the production
path, never assumed.
\end{remark}

\section{Draw $\times$ field size: the crowding interaction}

The geometric penalty is amplified when the field is crowded: a wide draw costs more in a full
field (more bodies to be trapped behind, fewer gaps to cross to the rail) than in a sparse one ---
the first \emph{interaction} moment of the draw mechanism. The production form scales the draw gain
by a field-size factor centred at twelve runners, clamped so the penalty never reverses in a small
field:
\begin{equation}
  g_{\text{draw}}^{\text{eff}} \;=\; g_{\text{draw}}\cdot
   \max\!\bigl(0,\ 1 + g_{\text{fs}}\,(N-12)\bigr),
  \label{eq:drawfs}
\end{equation}
with $\texttt{DRAW\_FS\_GAIN\_BY\_VENUE}$ (HV $0.25$, ST $0.0$). It converts only at HV (the
tightest track, where collision physics dominates: HV $+0.3$pp, flat at ST). It is explicitly an
\emph{interim stopgap} --- a scalar shadow of a genuinely two-dimensional geometry$\times$control
interaction (Chapter~\ref{ch:theory}) --- and is flagged as such pending the joint treatment.

\section{Draw as variance: loading the diffusion knob}

The deeper insight is that draw acts not only on the \emph{mean} (the arc-length drift penalty) but
on the \emph{variance}: a wide-drawn horse in a big field has a more uncertain trip --- it may get
across and save ground, or it may be hung wide and buried --- and that uncertainty is a property of
$\Ddiff$, not of the drift. This is the first deliberate use of the third knob. Each horse's
diffusion is inflated by a draw-derived ``trip risk'',
\begin{equation}
  \Ddiff_i \;=\; D_{\text{eff}}\,\bigl(1 + g_{\text{dvar}}\cdot
   \mathrm{bend\_fraction}(L)\cdot \mathrm{extremity}_i\bigr),
  \qquad
  \mathrm{extremity}_i = \frac{2\,\bigl|\mathrm{Dr}_i - \tfrac{N+1}{2}\bigr|}{N-1},
  \label{eq:drawvar}
\end{equation}
where $\mathrm{extremity}_i\in[0,1]$ is the field-normalised lateral extremity of the draw (a
rail or widest-out horse is extreme, a midfield draw is not), and the risk is gated by the bend
fraction (no bends, no trip risk). A larger $\Ddiff_i$ means a larger Wald variance
\eqref{eq:wald-var} and so more upsets from that horse. The gain
$\texttt{DRAW\_VAR\_GAIN\_BY\_VENUE}$ is ST $0.5$, HV $0.0$.

The ST-only split is \emph{predicted by theory}, a clean theory-confirms-data check. The winner is
the left-tail extreme of $N$ finishing-time draws; in a larger field the win requires a deeper tail
event, so per-horse variance has more leverage on \emph{who wins the larger $N$ is}. ST runs the
big fields, HV the small ones, so the variance knob should pay on ST and barely at all on HV ---
exactly the observed split (see Chapter~\ref{ch:theory}, the extreme-value corollary).

\section{Surface as diffusion}

Surface enters the variance knob directly: the all-weather (dirt) surface is more diffusive than
turf --- kickback and a less uniform surface make times less repeatable --- so its diffusion is
scaled by $\texttt{awt\_diff\_mult}=2.0$,
\begin{equation}
  D_{\text{eff}} \;=\; \texttt{diff\_coef}\cdot
   \begin{cases} 2.0 & \text{AWT}\\ 1.0 & \text{turf,}\end{cases}
   \qquad \texttt{diff\_coef}=170,
  \label{eq:awtdiff}
\end{equation}
which corrects a pronounced AWT over-confidence (raising the model's AWT win-probability dispersion
from an over-tight $12.7\%$ toward a calibrated $23.7\%$). The empirical distance-law study of
Chapter~\ref{ch:state} independently recovers an AWT diffusion roughly twice the turf value,
vindicating this multiplier from data.

\section{What geometry does and does not yet capture}

The terms above --- $\Leff$ on the drift, the crowding amplifier, the draw$\to\Ddiff$ variance, and
the surface diffusion --- are the \emph{mean-field shadows} of draw: per-horse scalars that capture
the field-averaged, first-moment consequences of geometry. They do \emph{not} capture the genuine
two-dimensional, many-body phenomenon in which the cost of a draw is a functional of the jockey's
control policy and of the style composition of the neighbouring horses. That is the
geometry$\times$control game of Chapter~\ref{ch:theory}, the single most promising piece of open
theory, for which the baked draw terms are explicit interim stopgaps and for which the
computer-vision trajectory programme is being built to supply the missing lateral coordinate.
